\SourcePage{232}{237}
\section{Radiation of nuclei}

For $\gamma$-radiation of nuclei, as a rule, the condition of smallness of the dimensions of the system (the nuclear radius $R$) compared with the wavelength of the photon is satisfied. However, the spacings between nuclear levels (and thereby the energy of the $\gamma$-quantum) are usually small compared with the energy per nucleon in the nucleus. Therefore the quantity $R/\lambda$ is not connected directly with the speed $v/c$ of nucleons in the nucleus and, generally speaking, is significantly smaller. Correspondingly

\SourcePage{233}{238}
the probability of $M1$-radiation is, as a rule, greater than that of $E$, $l + 1$ radiation (cf.\ the beginning of~\S\,50).

The general selection rules for the total angular momentum (``spin'') of the nucleus and for parity are the same as for radiation by any system. A characteristic feature of nuclear radiation is the prevalence of transitions of higher multipolarity. In contrast to atoms, whose radiation is usually electric-dipole, in nuclei at small energies such transitions are comparatively rare, being forbidden by the selection rules.

If a radiative transition of a nucleus can be regarded as a single-particle process---a change of state of one nucleon with the state of the nuclear ``core'' remaining unchanged---then the selection rules for the angular momentum of this nucleon are added. The accuracy of observance of such ``single-particle'' selection rules turns out, however, to be very low.

Specific for nuclei are the selection rules for the isotopic spin. Recall that the projection $T_3$ of the isotopic spin is determined already by the mass and number of the nucleus:
\[
T_3 = \tfrac{1}{2}(Z - N) = Z - \tfrac{A}{2}.
\]

For a given value of $T_3$ the absolute magnitude of the isotopic spin can have any values $T \geqslant |T_3|$. The selection rule for the number $T$ for radiative transitions arises in connection with the fact that the operators of the electric and magnetic moments of the nucleus, expressed through the operators of isotopic spins of the nucleons, represent sums of a scalar and the $x_3$-component of a vector in isotopic space (see~III, \S\,116). Therefore their matrix elements are non-zero only under the condition
\begin{equation}
T' - T = 0,\,\pm 1.
\tag{55.1}
\end{equation}
In itself this rule, however, does not impose any particular restrictions on transitions in light nuclei (for which only one can speak with sufficient accuracy about the conservation of isotopic spin); the fact is that among the lowest levels of these nuclei there are practically no levels with $T > 1$.

But for $E1$-transitions there is an additional rule in connection with the fact that for the electric dipole moment the isotopically scalar part drops out, and its operator reduces to the $x_3$-component of the isotopic vector (see~III, \S\,116). Therefore if $T_3 = 0$, the transitions with $\Delta T = 0$ are additionally forbidden. In other words, in nuclei with equal numbers of neutrons and protons ($N = Z$, $A = 2Z$) $E1$-transitions are possible only when
\begin{equation}
T' - T = \pm 1\qquad (T_3 = 0).
\tag{55.2}
\end{equation}
Of course, the accuracy of observance of this rule depends on the accuracy with which the isotopic spin of the nucleus is conserved.

The probability of $E1$-transitions in the nucleus is also influenced by the effect of recoil of the nuclear core in the motion of individual nucleons. This effect leads to the fact that in forming the dipole moment the protons participate with an effective charge $e(1 - Z/A)$ instead of $e$, and the neutrons with a charge $-eZ/A$ instead of 0 (see~

\SourcePage{234}{239}
III, \S\,118). The reduction of the effective charge of the proton leads to a certain suppression of the probability of $E1$-transitions.

The energy levels of non-spherical nuclei possess rotational structure. In connection with this there appears a specific rotational structure of the spectrum of $\gamma$-radiation for such nuclei. The symmetry of the field in which the nucleons move in a ``stationary'' non-spherical (axially symmetric) nucleus coincides with the symmetry of the field in which the electrons move in a ``stationary'' diatomic molecule of identical atoms (point group $\mathrm{C}_{\infty h}$). Therefore the symmetry properties of levels of the non-spherical nucleus (and with them the selection rules for matrix elements) are analogous to the symmetry properties of levels of a diatomic molecule (see~III, \S\,119). In particular, as for a diatomic molecule of identical atoms, electric-dipole transitions are forbidden within one and the same rotational band (i.e.\ without change of the internal state of the nucleus)---cf.~\S\,54. These transitions are therefore realised as $E2$- or $M1$-transitions. In the first case the total angular momentum $J$ of the nucleus can change by 2 or 1, and in the second by~1.

According to~(46.9) the probability of the quadrupole transition, summed over the final values of the projection $M'$ of the total angular momentum of the nucleus, is
\[
w_{E2} = \frac{\omega^5}{15\hbar c^5}\sum_{M'}|\langle J'\Omega M'|Q^{(\mathscr{E})}_{2,-m}|J\Omega M\rangle|^2
\]
($J$ is the total angular momentum of the nucleus; $\Omega$ is its projection onto the axis of the nucleus; $m = M - M'$). With the help of formula~(110.8) (see~III) this sum can be expressed through the squares of the given quantities---the diagonal (with respect to the internal state of the nucleus) quadrupole moments of the transition $\bar{Q}_{2\lambda}$, defined in the coordinate system $\xi\eta\zeta$ associated with the nucleus. Here $\lambda = \Omega - \Omega'$, so that in the given case ($\Omega' = \Omega$) only the component $\bar{Q}_{20}$ appears. By definition the quantity $e\bar{Q}_0$ is simply called the quadrupole moment of the nucleus:
\[
eQ_0 = e\int\rho_{ii}(2\zeta^2 - \xi^2 - \eta^2)\,d\xi\,d\eta\,d\zeta = -2e(\bar{Q}_{20})_{ii}.
\]
We therefore obtain
\begin{equation}
w_{E2}(\Omega J \to \Omega J') = \frac{\omega^5}{60\hbar c^5}\,Q^2_0(2J'+1)\begin{pmatrix} J' & 2 & J \\ -\Omega & 0 & \Omega\end{pmatrix}^{\!2}.
\tag{55.3}
\end{equation}

\SourcePage{235}{240}
Concerning these formulae, however, the following remark should be made. In them the matrix elements computed with wave functions of the form
\[
\psi_{J\Omega M} = \mathrm{const}\cdot\chi_\Omega\,D^{(J)}_{\Omega M}(\mathbf{n})
\]
($\chi$ is the wave function of the internal state of the nucleus) have been used. These functions correspond to definite (in magnitude and sign) values of the projection of the angular momentum onto the axis $\zeta$. In nuclei, however, we are dealing with states possessing only definite parity and magnitude of the projection of the angular momentum (the latter quantity is usually and will be understood under $\Omega$). Therefore as initial and final wave functions one should take the combinations
\[
\frac{1}{\sqrt{2}}\,(\psi_{J\Omega M} \pm \psi_{J,-\Omega,M}),\qquad \frac{1}{\sqrt{2}}\,(\psi_{J'\Omega M'} \pm \psi_{J',-\Omega,M'}).
\]
The products of first and second terms will give the same value of the matrix element of the quadrupole transition. The ``cross'' terms, however, will lead to non-vanishing integrals if $2\Omega \leqslant 2$\,\footnote{For the matrix elements of the $l$-pole transition the sub-integral expressions will contain products of the form
\[
D^{(J')}_{-J,\Omega M'}\,D^{(l)}_{q'q}\,D^{(J)}_{\Omega M}.
\]
The integral over angles will be non-zero only when $q' = -2\Omega$, while $q'$ runs only from $-l$ to $l$; therefore it is necessary that $2\Omega \leqslant l$.}. Therefore formula~(55.3), strictly speaking, is inapplicable for $\Omega = \tfrac{1}{2}$, 1; in these cases an additional term appears in the probability of the transition, not expressible through the mean value of the quadrupole moment\,\footnote{In practice this correction gives a substantial correction only for $\Omega = \tfrac{1}{2}$, when the connection between the rotation and the internal state of the nucleus is especially strong (see~III, \S\,119).}.

Analogously to the derivation of formula~(55.3) for the probability of the $M1$-transition we obtain the formula
\begin{equation}
w_{M1}(\Omega J \to \Omega,\,J-1) = \frac{4\omega^3}{3\hbar c^3}\,\mu^2(2J-1)\begin{pmatrix} J-1 & 1 & J \\ -\Omega & 0 & \Omega\end{pmatrix}^{\!2} = \frac{4\omega^3}{3\hbar c^3}\,\mu^2\,\frac{J^2 - \Omega^2}{J(2J+1)},
\tag{55.4}
\end{equation}
where $\mu$ is the magnetic moment of the nucleus (this formula is inapplicable for $\Omega = \tfrac{1}{2}$).
